\documentclass[10pt,twocolumn]{article}

\usepackage[margin=0.75in]{geometry}
\usepackage{amsmath,amssymb,amsthm}
\usepackage[numbers]{natbib}
\usepackage{microtype}
\usepackage{hyperref}
\usepackage{booktabs}

\newtheorem{theorem}{Theorem}
\newtheorem{definition}{Definition}

\newcommand{\Csmooth}{\bar{C}}
\newcommand{\Craw}{C}
\newcommand{\tasktype}{\tau}

\title{Backpressure Economics:\\
Congestion Control for Agent Payment Flows}
\author{Pura Protocol}
\date{\today}

\begin{document}
\maketitle
\thispagestyle{empty}

% ═══════════════════════════════════════════════════
% PAGE 1 — Problem + competitors
% ═══════════════════════════════════════════════════

\section{Problem}

Streaming payment protocols let AI agents pay each other continuously, but none
of them handle congestion. When a downstream provider is at capacity, payments
pile up with no mechanism to reroute, buffer, or throttle. Data networks solved
this decades ago with TCP congestion control. Payment networks have not.

The result: the best providers get flooded while mediocre ones sit idle. Callers
have no signal about who has spare capacity. Money flows in, but verified work
does not always come out.

\subsection{Existing approaches}

\textbf{x402} (Coinbase, 2025) wraps HTTP 402 responses with crypto payment
headers. No capacity signal, no routing, no flow control. The client pays and
hopes the server has room.

\textbf{Tempo MPP} adds a temporary multi-path payment layer to Lightning.
Splits help with channel liquidity but provide zero information about
downstream service capacity.

\textbf{AP2/TAP} (The Graph) handles micropayments for indexed blockchain data.
Fixed pricing, no congestion response, no capacity attestation.

\textbf{Load balancers} (Nginx, HAProxy, cloud LBs) read server-side health
checks and distribute requests. The capacity signal stays behind the firewall.
No on-chain verification. No payment integration. No cross-org trust.

None of these combine three things: a public capacity signal, cryptographic
completion verification, and congestion-responsive routing.

% ═══════════════════════════════════════════════════
% PAGE 2 — Mechanism
% ═══════════════════════════════════════════════════

\section{Mechanism}

Backpressure Economics (BPE) adapts the Tassiulas--Ephremides backpressure
routing algorithm~\cite{tassiulas1992stability} from communication networks
to monetary flows. Five on-chain primitives do the work.

\subsection{Capacity signal (EWMA)}

Each provider declares capacity $\Craw(k,t)$. A commit-reveal scheme prevents
gaming. The protocol smooths raw signals with an exponentially weighted moving
average:
\begin{equation}
\Csmooth(k,t) = \alpha \cdot \Craw(k,t) + (1-\alpha) \cdot \Csmooth(k,t-1)
\end{equation}
with $\alpha = 0.3$, matching Jacobson's TCP RTT estimator~\cite{jacobson1988congestion}.
Smoothing damps oscillation while tracking real capacity shifts within a few epochs.

\subsection{Virtual queue}

For each provider $k$, a virtual queue tracks unprocessed payment backlog:
\begin{equation}
Q(k,t\!+\!1) = \max\{Q(k,t) + F(k,t) - \Csmooth(k,t),\; 0\}
\end{equation}
where $F(k,t)$ is payment flow routed to $k$ at time $t$. The queue is virtual
--- no tokens are locked --- but it records how overloaded each provider is.

\subsection{Max-weight routing}

The total incoming payment flow $\Lambda(t)$ is split across providers
proportionally to their capacity-weighted differential backlog:
\begin{equation}
F(k,t) = \frac{W(k,t)^{+} \cdot \Csmooth(k,t)}
{\sum_{k'} W(k',t)^{+} \cdot \Csmooth(k',t)} \cdot \Lambda(t)
\end{equation}
where $W(k,t) = Q_{\text{source}}(t) - Q(k,t)$ is the differential backlog and
$x^{+} = \max(x,0)$. In steady state (all backlogs similar), this reduces to
proportional-to-capacity allocation --- the mechanism implemented in Superfluid's
General Distribution Agreement.

\subsection{Overflow buffer}

Money cannot be dropped like a packet. When aggregate demand exceeds aggregate
capacity, excess payment enters an on-chain escrow buffer $B(t)$. The buffer
drains proportionally as capacity frees up. Demurrage ($\delta = 2.5\%$/yr)
decays idle buffer funds, preventing permanent lockup.

\subsection{Dynamic pricing}

Price rises with congestion, following an EIP-1559-style adjustment:
\begin{equation}
p(k,t) = \beta(t) \cdot \left(1 + \gamma \cdot \frac{Q(k,t)}{\Csmooth(k,t)}\right)
\end{equation}
The base fee $\beta$ increases 12.5\% per epoch when demand exceeds capacity and
decreases 12.5\% otherwise, converging to the shadow price from Kelly's
proportional fairness framework~\cite{kelly1998rate}.

% ═══════════════════════════════════════════════════
% PAGE 3 — Guarantee
% ═══════════════════════════════════════════════════

\section{Throughput guarantee}

The core result is that BPE serves every demand vector that the provider
network can physically handle, with bounded overflow.

\begin{theorem}[Throughput optimality]
\label{thm:main}
For any arrival rate vector $\boldsymbol{\lambda}$ strictly inside the capacity
region (i.e., $\boldsymbol{\lambda} \in (1-\epsilon)\Lambda^{*}$ for some
$\epsilon > 0$), max-weight payment routing stabilizes all virtual queues:
\begin{equation}
\limsup_{T \to \infty} \frac{1}{T} \sum_{t=0}^{T-1} \sum_{k}
\mathbb{E}[Q(k,t)] \leq \frac{D}{\epsilon}
\end{equation}
where $D$ depends on maximum arrival and capacity rates.
\end{theorem}

\textit{Proof sketch.} Define the Lyapunov function $L(t) = \sum_k Q(k,t)^2$.
From queue dynamics, the one-step drift satisfies
$\Delta(t) \leq D + 2\sum_k Q(k,t) \cdot (F(k,t) - \Csmooth(k,t))$.
Max-weight minimizes the second term over all feasible allocations. Since
$\boldsymbol{\lambda}$ is strictly interior, a stationary policy exists with
slack $\epsilon' > 0$. Max-weight does at least as well, yielding
$\Delta(t) \leq D - \epsilon' \sum_k Q(k,t)$. Foster--Lyapunov gives
positive recurrence; telescoping the drift gives the bound.
See Neely~\cite{neely2010stochastic}, Ch.~4 for the full argument.

The overflow buffer inherits the same bound via an augmented Lyapunov function
$\tilde{L}(t) = L(t) + \sum_\tau B(\tau,t)^2$.

\textit{What this means in practice:} if providers collectively have enough
capacity to handle the incoming payment load, BPE will find the allocation.
Requests land on providers with spare room. Payments flow where work actually
gets done. The buffer stays small.

% ═══════════════════════════════════════════════════
% PAGE 4 — Implementation + results
% ═══════════════════════════════════════════════════

\section{Implementation}

BPE is implemented as the Pura protocol: 32 Solidity contracts deployed on
Base Sepolia, tested with 319 passing tests across 32 suites.

\subsection{Contracts}

The core contracts: \texttt{CapacityRegistry} (EWMA attestations),
\texttt{CompletionLedger} (dual-signed receipt verification),
\texttt{BackpressurePool} (Superfluid GDA wrapper with max-weight routing),
\texttt{OverflowEscrow} (buffer with demurrage), and
\texttt{QueueLengthPricing} (EIP-1559-style dynamic pricing).

Extensions include \texttt{EconomyFactory} (one-tx pool deployment),
\texttt{NestedPool} (pipeline composition with upstream congestion propagation),
\texttt{QualityOracle}, and a thermodynamic layer (\texttt{TemperatureOracle},
\texttt{VirialMonitor}) that maps attestation variance to system temperature
and triggers phase-transition circuit breakers.

\subsection{Gateway}

The Pura gateway is an OpenAI-compatible HTTP endpoint (\texttt{api.pura.xyz})
that routes inference requests across providers using on-chain capacity weights.
Drop-in replacement: change \texttt{baseURL} to \texttt{https://api.pura.xyz/v1}
in any OpenAI SDK client. Every completion is recorded on-chain with the
provider, model, and token count.

\subsection{Simulation results}

Agent-based simulations over 1,000-step horizons with market-phase variation
(bull, bear, shock, recovery):

\begin{center}
\begin{tabular}{@{}lcc@{}}
\toprule
Metric & Round-robin & BPE \\
\midrule
Allocation efficiency & 93.5\% & 95.7\% \\
Buffer stall rate & 73\% & $<$9\% \\
Shock recovery & --- & 50 steps \\
Sybil profitability & --- & negative $\forall n > 1$ \\
\bottomrule
\end{tabular}
\end{center}

Buffer stall rate drops from 73\% to under 9\% with overflow escrow sized at
one period of peak demand. Sybil analysis confirms stake fragmentation yields
strictly negative profit for all split counts $n > 1$ due to the concave-sqrt
weighting function.

\subsection{Status}

Testnet contracts are live. The gateway routes real inference requests. SDK
(\texttt{@puraxyz/sdk}) provides 23 action modules covering all contract
interactions. Mainnet deployment is next, pending audit.

Full paper, source code, and simulation: \url{https://pura.xyz}

\bibliographystyle{plainnat}
\bibliography{../bpe}

\end{document}
